\begin{lemma}[Exact lifting for covariant polynomial tensors]
\label{c:covariant-polynomial-lifting}
\status{Proved}
Let $V_0=\Q^{u|v}$, let $L_0$ be an even line, and set
$G=\operatorname{GL}(V_0)\times\mathbb G_m$.  Fix $q\geq1$ and put
\[
 W=\bigoplus_{m=1}^{q}(V_0^\vee)^{\otimes m}\otimes L_0,
 \qquad R=\mathcal O(W)
 =\operatorname{Sym}_{\mathrm{super}}
       \left(\bigoplus_{m=1}^{q}V_0^{\otimes m}\otimes L_0^{-1}\right).
\]
Grade $R$ by total $V_0$-degree $r$ and number $d$ of form factors.
Every $R_{r,d}$ is finite and is a direct summand of a finite sum
of copies of $V_0^{\otimes r}\otimes L_0^{-d}$.
If $I\subset R$ is a $G$-stable ideal and $E\subset I$ is a finite
homogeneous $G$-submodule generating $I$, then
\begin{equation}
 \Hom_G\bigl(V_0^{\otimes r}L_0^{-d},(E\otimes R)_{r,d}\bigr)
 \longrightarrow
 \Hom_G\bigl(V_0^{\otimes r}L_0^{-d},I_{r,d}\bigr)
 \label{eq:c-polynomial-lift}
\end{equation}
is surjective.  The same conclusion holds after a finite direct
sum in the source.  All the splittings and lifts can be found by
rational linear algebra.
\end{lemma}

\begin{proof}
The degree-$r$ polynomial representations of
$\operatorname{GL}(u|v)$ are modules for its Schur superalgebra.
The coefficient-coalgebra description and the explicit
coefficient-basis proof in
\citet[section~5, (5.2), Lemma~5.1 and Theorem~5.2,
author pp.~13--15]{c_brundankujawa2003}
identify that algebra, over $\Q$, with
\[
 S(u|v,r)=\End_{\Q[\mathfrak S_r]}(V_0^{\otimes r}),
\]
where the symmetric group acts by superpermutations.
Maschke's theorem makes this a semisimple algebra: decompose the
finite $\Q[\mathfrak S_r]$-module into its simple isotypical
parts and take the product of their matrix endomorphism algebras.
This is also the characteristic-zero argument in
\citet[Theorem~1.3.1(a)(i), proof of Theorem~4.2.6(a),
author pp.~3 and~12]{c_hemmerkujawanakano2006}.
An ungraded splitting can be made even by replacing a map $f$
by $(f+\Pi f\Pi)/2$, where $\Pi$ denotes intrinsic parity on the
source and target.  Conjugation by $\Pi$ preserves the module
intertwiner equations, so this averaging still gives a splitting.
Thus the category with even morphisms is semisimple as well.
Tensoring by a fixed character of the independent last torus does
not change this conclusion.

The degree-$d$ supersymmetric power is the image of the averaging
idempotent on the $d$-fold tensor power of its generators.  After
fixing total degree $r$, the latter is a finite sum of copies of
$V_0^{\otimes r}L_0^{-d}$.  Finiteness follows also directly from
$d\leq r\leq qd$, except for $R_{0,0}=\Q$.
The two central scalar actions split every stable ideal into
these homogeneous components.  The multiplication surjection
$(E\otimes R)_{r,d}\twoheadrightarrow I_{r,d}$ is therefore a
surjection of finite polynomial representations of the same
degree and Tate character.  Its source, target and the source of
the displayed Hom functor belong to the semisimple category just
proved.  It splits, proving the assertion.  Intertwining and
splitting impose finite rational linear equations.  We have used
exactness in the covariant polynomial category, not exactness of
invariants on arbitrary mixed rational superrepresentations.
\end{proof}

\begin{theorem}[A finite criterion without a sign hypothesis]
\label{c:finite-unconditional-criterion}
\status{Proved; finite existence is unconditional, and effectiveness has the stated input condition}
Let $X/k$ be smooth projective and connected, of dimension $n>0$
over an algebraically closed field, let $H$ be a Weil cohomology
with coefficient field $K$, and let
$o\in\CH^n(X)_\Q$ have degree one.  Set
\[
 a=\dim_K H^{\mathrm{even}}(X),\qquad
 b=\dim_K H^{\mathrm{odd}}(X),\qquad
 N=a(b+1),\qquad \lambda=((b+1)^a).
\]
In Chow motives put $A=h(X)$, $L=\Q(-n)$,
\[
 \epsilon(x)=\operatorname{tr}(o\cdot x),\qquad
 \pi=\id_A-\eta\epsilon,\qquad V=(A,\pi),
 \qquad
 \kappa_m=\operatorname{tr}\circ\mu^{(m)}\circ\pi^{\otimes m}
       :V^{\otimes m}\longrightarrow L.
\]
Here $\eta$ is the unit and $\mu^{(m)}$ is iterated cup product;
the displayed notation for $\epsilon$ includes its usual Tate
cancellation.  Since $\epsilon\eta=1$, $\pi$ is an actual Chow
projector, independently of any standard conjecture.  There is
a finite list of rational expressions $P_1,\ldots,P_s$ in the
forms $\kappa_1,\ldots,\kappa_{2n}$, using tensor products,
permutations and rational linear combinations, such that the
following are equivalent:
\begin{enumerate}[label=(\Alph*)]
\item All rational cohomological relations among the specified
      point and partial-diagonal generators on every power of $X$
      hold in rational Chow.
\item The following finite list of Chow identities holds:
      \begin{enumerate}[label=(\roman*)]
      \item $c_n(T_X)=(a-b)o$ and $\Delta_*o=o\boxtimes o$;
      \item $\kappa_{2n+1}=0$;
      \item the rectangular Young symmetrizer
            $y_\lambda:V^{\otimes N}\to V^{\otimes N}$ is zero;
      \item $P_\nu=0$ for $1\leq\nu\leq s$.
      \end{enumerate}
\end{enumerate}
Each identity in (ii)--(iv), after Poincar\'e-duality bending,
is a finite rational combination of cycles in the requested
point-and-diagonal rings.  The list can be constructed effectively
from a rational cohomology table, and more generally from a finite
presentation of the $\Q$-algebra generated by its structure
constants.  With an arbitrary unspecified Weil coefficient
field, finite existence is asserted, without an algorithm for
testing its constants.  No decidability assertion for rational
equivalence is made.  No algebraic parity projector is assumed.
More explicitly, the bent form $\kappa_m$ is the modified diagonal
\begin{equation}
 \Gamma_m(X,o)=
 \sum_{\varnothing\ne J\subseteq\{1,\ldots,m\}}
       (-1)^{m-|J|}\Delta_J\prod_{i\notin J}p_i^*o
       \ \in\CH^{n(m-1)}(X^m)_\Q,
 \label{eq:c-explicit-higher-modified-diagonal}
\end{equation}
where a singleton diagonal means the unit on that coordinate.
Thus condition (ii) is the explicit identity
$\Gamma_{2n+1}(X,o)=0$.
\end{theorem}

\begin{proof}
We first normalize the Frobenius algebra.  Poincar\'e duality
identifies the equality $\Delta_*o=o\boxtimes o$ with
\[
 \epsilon\mu=\epsilon\otimes\epsilon.
\]
Indeed the two sides pair with $x\otimes y$ as, respectively,
$\operatorname{tr}(oxy)$ and
$\operatorname{tr}(ox)\operatorname{tr}(oy)$.  This calculation
is an identity of Chow correspondences, since the pairing and
its inverse are the diagonal evaluation and coevaluation.
Consequently $V=\ker\epsilon$ is an algebra ideal, with a
restricted multiplication $\mu_V$.  For $m\geq1$ we have
\begin{equation}
 \kappa_{m+1}
   =\kappa_m\circ(\mu_V\otimes\id_V^{\otimes(m-1)}).
 \label{eq:c-high-form-recursion}
\end{equation}
Thus (ii) kills all $\kappa_m$ with $m>2n$ in Chow.  In
cohomology $H(V)=H^{>0}(X)$, so these same forms vanish by
degree for every $m>2n$.  This explains why a single additional
modified-diagonal identity suffices to truncate the forms.

Here is the finite construction.  Let $V_0=\Q^{a-1|b}$ and let
$L_0$ be a formal even line.  We do not assert that $V_0$ is a
rational form of the given cohomology algebra.  Apply
Lemma~\ref{c:covariant-polynomial-lifting} with $q=2n$ and write
$G=\operatorname{GL}(V_0)\times\mathbb G_m$.
The cohomological forms give an even $K$-point
$t=(\kappa_1,\ldots,\kappa_{2n})$ of $W$.  Define
\begin{equation}
 I_\Q=\ker\left(R\longrightarrow\mathcal O(G)\otimes_\Q K\right),
 \qquad f\longmapsto\bigl(g\longmapsto f(g\cdot t)\bigr).
 \label{eq:c-pure-form-orbit}
\end{equation}
This is the kernel of a rational comodule map, with the regular
action on $\mathcal O(G)$ and trivial action on $K$.  It is
therefore a $G$-stable ideal.  We use the schematic super-orbit:
odd directions and nilpotents are retained.  Taking only the
ordinary-point orbit would be a different and insufficient
construction.  The two central scalar actions make $I_\Q$
homogeneous in $(r,d)$.

The supercommutative algebra $R$ is a finite module over the
polynomial algebra on its even generators, and hence is
noetherian.  Choose finite ideal generators of $I_\Q$ and close
their homogeneous span under $G$.  This yields a finite
homogeneous stable submodule $E\subset I_\Q$ that still generates
the ideal.  Each $E_{r,d}$ lies in the finite semisimple module
$R_{r,d}$ and hence is a split summand of a finite sum of copies
of $V_0^{\otimes r}L_0^{-d}$.  Choose a finite such sum $T$ and
maps
\[
 s:E\longrightarrow T,\qquad p:T\longrightarrow E,
 \qquad ps=\id_E.
\]
For each tensor summand $T_\nu=V_0^{\otimes r_\nu}L_0^{-d_\nu}$,
the composite
\begin{equation}
 T_\nu\longrightarrow T\xrightarrow{p}E\hookrightarrow R
 \label{eq:c-finite-form-relations}
\end{equation}
is a polynomial multilinear form expression.  After evaluation
at $t$ and cancellation of the indicated line twist it is a map
$V_0^{\otimes r_\nu}\otimes K\to L_0^{d_\nu}\otimes K$.
Its Chow expression is $P_\nu$.

We justify both the expression and every equality among such
expressions.  In each bidegree the symmetric-power averaging
idempotents split $R_{r,d}$ into a finite sum of pure tensor
powers with its fixed line character.  The first fundamental
theorem in \citet[Theorem~2.1]{c_lehrerzhang2021super} says that
the intertwiners between equal pure powers are generated by
permutations.  Hence \eqref{eq:c-finite-form-relations} is a rational combination of products
of the universal forms, evaluated on disjoint groups of inputs.
The second fundamental theorem, their Theorem~2.3, says that all
relations between those permutation maps are generated by the
Young symmetrizer with $(a-1)+1=a$ rows and $b+1$ columns.
The central scalar excludes unequal pure degrees.  Inserting
the averaging idempotents merely gives direct summands of this
same description.  These statements descend from $\C$ to $\Q$
because the maps and their kernel equations are rational.
The fixed Tate line is even and invertible, and its usual
one-dimensional identities already hold in Chow motives.
Thus (iii) makes every representation-theoretic equality used
in this construction a valid Chow tensor identity.

Now let a rational universal form
\[
 F:V_0^{\otimes r}L_0^{-d}\longrightarrow R
\]
be equivariant and satisfy $F(t)=0$.  Equivariance on all
superalgebra-valued points gives $F(g\cdot t)=gF(t)=0$;
equivalently its image is contained in $I_\Q$.  Conversely,
membership in $I_\Q$ implies zero evaluation at the identity.
The multiplication map $E\otimes R\twoheadrightarrow I_\Q$
and Lemma~\ref{c:covariant-polynomial-lifting} give a lift
\[
 \widetilde F:V_0^{\otimes r}L_0^{-d}
                 \longrightarrow(E\otimes R)_{r,d}.
\]
Compose with $s\otimes\id$ to obtain a lift with target
$(T\otimes R)_{r,d}$.  Expanding this lift with the pure-tensor
first fundamental theorem expresses $F$ as a sum of expressions
that apply one of the maps \eqref{eq:c-finite-form-relations} to some inputs and other
universal forms to the remaining inputs.  All intervening
objects are covariant polynomial representations of total
$V_0$-degree $r$ and line character $-d$; no exactness assertion
for an arbitrary mixed invariant functor has entered.
The equality is valid as a formal tensor identity modulo the
rectangular relation by the preceding second fundamental
theorem.  Conditions (iii)--(iv) therefore imply
$F(\kappa_1,\ldots,\kappa_{2n})=0$ in Chow.  There is no
projective auxiliary tensor to cancel and no use of $C^+(X)$.

We next return from forms to the requested cycles.  By (i) and
the graph product rule, every monomial in the generators reduces
to a scalar times a pointed-partition cycle.  A homogeneous
relation between these cycles, with $k$ coordinates and $d$
unpointed blocks, bends to a map
$A^{\otimes k}\to L^d$.  For an unpointed block $B$ its factor
is $\operatorname{tr}(\prod_{i\in B}x_i)$; for a pinned
coordinate it is $\epsilon(x_i)$.  Split every input using the
actual decomposition $A=\mathbf1\oplus V$.  A pinned $V$ input
annihilates its term.  Unit inputs in an unpointed block are
deleted; if no input remains, its value is
$\operatorname{tr}(1)=0$, since $n>0$.  Every other block is one
of the forms $\kappa_m$.  Condition (ii) removes blocks with
$m>2n$.  Each remaining restriction is exactly a rational form
$F$ of the preceding paragraph.  If the original relation is
cohomologically zero, each such restriction vanishes at $t$,
and hence vanishes in Chow.  The direct-sum decomposition of
$A^{\otimes k}$ and Poincar\'e duality prove the original cycle
relation.  This proves (B)$\Rightarrow$(A).

For necessity, the cohomological relations
$\Delta_*o-o\boxtimes o=0$ and
$\Delta_X^2-(a-b)o\boxtimes o=0$ belong to the requested ring
on $X^2$.  Pushing the second along a projection and using
self-intersection gives $c_n(T_X)=(a-b)o$.  The form
$\kappa_{2n+1}$ is cohomologically zero by positive degrees.
Its bent Chow cycle is obtained from the small diagonal by
inserting $\pi$ in each input.  Expanding
$\pi=\id-\eta\epsilon$ gives pointed partial-diagonal cycles.
The same expansion shows that $y_\lambda(V)$, viewed as a
correspondence on $X^N$ and hence a cycle on $X^{2N}$, is a
rational combination of products of pair diagonals and point
classes.  It is cohomologically zero by the rectangular theorem.
Finally every $P_\nu$ is a product of the forms just described,
with permuted inputs, and is cohomologically zero because
$E\subset I_\Q$.  Bending gives the claimed cycles in the
requested rings.  Thus (A) implies (ii)--(iv) as well as (i).
For the displayed formula for $\Gamma_m$, the adjoint of
$\pi=1-\eta\epsilon$ under Poincar\'e duality is
$1-o\operatorname{tr}$.  Expanding its action on the small
diagonal gives the terms indexed by nonempty $J$ in
\eqref{eq:c-explicit-higher-modified-diagonal}.  The omitted empty
term integrates every coordinate and is
$\operatorname{tr}(1)=0$.

For effectiveness, the orbit map is computed using two even
matrix blocks, odd off-diagonal entries, inverses of the two
even determinants, and the last torus coordinate and its inverse.
If the coordinates of $t$ generate a presented algebra
$B=\Q[c_1,\ldots,c_h]$, compute the kernel into
$\mathcal O(G)\otimes B$ by orbit-graph elimination, including
the defining equations of $B$.  This is finite-module
polynomial elimination over the even variables: the odd
variables have a finite exterior basis, with known signed
multiplication.  Noetherian module Gr\"obner elimination
terminates.  The kernel agrees with \eqref{eq:c-pure-form-orbit}
because $B\hookrightarrow K$.  Stable spans, the splittings,
permutation expressions and all lifts then use finite rational
linear algebra.  Without such a presentation only the
noetherian existence argument is used.  This completes the proof.
\end{proof}

\begin{corollary}[A finite power detects all the relations]
\label{c:finite-power-cutoff}
\status{Proved; the cut-off is effective with the same coefficient input}
In Theorem~\ref{c:finite-unconditional-criterion}, let $r_\nu$
be the input arity of $P_\nu$ and put
\[
 B=\max\{2,\ 2n+1,\ 2a(b+1),\ r_1,\ldots,r_s\}.
\]
Then all rational cohomological relations hold in Chow on every
power of $X$ if and only if the cycle map is injective on
$T_B^{\CH}(X,o)$.  In particular such a finite cut-off exists
without $C^+(X)$ or a rational form of the Weil cohomology.
Its effective construction has exactly the coefficient-data
qualification in the theorem.
\end{corollary}

\begin{proof}
All finite tests in the theorem are bent cycles on powers of
arity at most $B$.  The two normalization tests are detected by
the indicated relations on $X^2$, the high form by $X^{2n+1}$,
the rectangle by $X^{2N}$, and $P_\nu$ by $X^{r_\nu}$.
Injectivity on $T_B^{\CH}$ implies injectivity on every
$T_k^{\CH}$ with $k\leq B$: pull back along $p:X^B\to X^k$,
and retract that pullback by
\[
 z\longmapsto p_*\left(z\prod_{i=k+1}^{B}p_i^*o\right).
\]
The composite on $\CH^*(X^k)_\Q$ is the identity because
$\deg o=1$.  Thus injectivity on the single specified power
implies all the finite tests, and the theorem applies.  The
converse is immediate.  The cut-off depends only on the
cohomological form data and the choices in its finite
construction, not on an unproved Chow vanishing.
\end{proof}

\begin{remark}[What the finite criterion settles]
\label{c:finite-unconditional-scope}
\status{Proved scope statement; the actual general Chow vanishing remains unresolved}
Theorem~\ref{c:finite-unconditional-criterion} removes the sign
hypothesis from the finite criterion, including for an arbitrary
Weil coefficient field.  It supplies finite generators for the
all-power obstruction and does not assert that these cycles
vanish for every variety.  It also does not make rational
equivalence decidable.  Its list is not claimed minimal; in
particular the more precise K3 obstruction identified above is
unchanged.  The explicitly requested all-power K3 assertion and
the unrestricted evaluation of the criterion for every $(X,o)$
therefore remain separate unresolved parts of C.1(c).
If $n=0$, connectedness over the algebraically closed field gives
$X=\operatorname{Spec}(k)$, and every relation holds immediately.
\end{remark}
